\odiseachapter{circe}{La isla de la hechicera}\label{cap:circe}

¿Y si la gran pasión de Odiseo no hubiera sido Penélope, sino Circe? ¿Y si su viaje a Ítaca no hubiera sido un regreso al hogar, sino la fuga de una cárcel de amor?

Desde luego, no es eso lo que nos cuenta Nolan. Su episodio da muchos detalles bien observados, entre ellos la isla llena de falsos animales, o la transformación de los compañeros de Ulises en cerdos aún antes de perder la forma humana, reflejada por el ansia con la que devoran el guisado. Pero su retrato de Circe sigue el mismo patrón que en el caso de Polifemo: prefiere presentarnos un personaje simple, casi plano, cuyas motivaciones son una mezcla de miedo y odio genérico a los hombres. A cambio de la superficialidad psicológica, nuestro fílmico director nos obsequia con toda una exhibición de ingeniería plástica, en la que la bruja transforma, \emph{a mano} y sudando lo suyo, a los hombres en cerdos.

Nada más. La relación entre Circe y Ulises es una simple transacción comercial ---mis hombres por tu hermana transformada en urraca---, adornada con unas cuantas frases de compromiso. Nolan no nos explica quién es Circe ---a pesar de que se trata del personaje más sobrenatural del relato, en cuanto que su magia es explícita, a diferencia de la de Calipso, por ejemplo---, ni nos informa de su agenda. La visita a la isla es breve y el viaje a los infiernos que sigue se muestra como un episodio independiente.

Veamos qué nos cuenta Homero. A Eea llega Odiseo con una sola nave, la única que escapó de los lestrigones:

\begin{versocita}
Y así llegamos a la isla de Eea; en ella habitaba\\
Circe beltrenza, la fascinadora deidad vocihumana,\\
que era de Eetes mientesfunestas verísima hermana;\\
los dos nacidos del Sol, que la luz a mortales regala,\\
y ambos por madre venidos de Perse, de Océano alhaja.
\end{versocita}

Fíjate, observadora lectora, detallista lector, cómo nos presenta Homero a la bruja beltrenza. La palabra clave que la define es \gr{δεινή} (deinḗ), literalmente temible o terrible, pero también imponente y venerable, términos que Tobal condensa en «fascinadora». También es importante «vocihumana», traducción afortunada de \gr{αὐδήεσσα} (audḗessa) que asegura que Circe es capaz de hablar con los hombres (y de embaucarlos, como veremos) en su propia lengua. Y, por supuesto, su linaje es formidable, pues es hija del Sol y de la deidad marina Perse (o Perseis en Hesíodo), es decir, una diosa de pleno derecho.

Recién llegados, Ulises se va a explorar, distingue el palacio de Circe ---nada de cabaña--- en lo alto de una peña y luego caza un ciervo con el que él y sus hombres se dan un festín en la playa. Al día siguiente los anima a ir a explorar, pero la tripulación, después de habérselas visto con cíclopes y lestrigones, no está para muchas aventuras. Finalmente forman dos grupos. Uno al mando del propio Odiseo y el otro capitaneado por su lugarteniente Euríloco. Echan a suertes a quién le toca explorar y la suerte cae en este último, que parte con muy pocas ganas:

\begin{versocita}
Conque marchó, y juntamente con él veintidós compañeros\\
llorando; atrás nos dejaban allí a nosotros gimiendo.
\end{versocita}

Viene ahora la descripción de la fauna que rodea el palacio:

\begin{versocita}
Dieron por una honda hoz con la casa de Circe, unos lienzos\\
hechos de piedra pulida, y, en torno, un lugar bien abierto.\\
Y por allí dormitaban leones y lobos serreños\\
que ella tenía hechizados por obra de algún mal veneno.\\
Pues, sin amago de echarse a mis hombres, en pie se pusieron\\
para halagarlos meneando sus colas al paso de ellos.\\
Y como, cuando al volver de una cena su amo, los perros\\
mueven el rabo, pues siempre él les trae golosina, contentos,\\
así los leones allí y los lobos pezuñirrecios\\
fiestas hicieron; miraban aquellos los mostros con miedo.
\end{versocita}

Nolan nos muestra los animales, sugiriendo que se trata de humanos convertidos, como también lo es el ciervo que caza Ulises. Homero no lo deja claro ---desde luego el ciervo no se convierte en humano---, pero el detalle no es importante. Las pacíficas fieras revelan que los aqueos se están acercando a la morada de una hechicera peligrosa. Pero los siguientes versos dan un giro al argumento.

\begin{versocita}
Ya en el portón de la trencigalana deidad, se tuvieron,\\
y a Circe oían cantar ---era hermosa su voz--- allí dentro,\\
puesta al telar a una tela inmortal ---de una diosa era empleo---\\
fina y graciosa y brillante, que así son las obras del cielo.
\end{versocita}

No hay en Homero una granjera que vuelva a su choza acarreando sus bártulos, sino una deidad de hermosa voz ---el detalle me impactó tanto que mi propia versión de Circe se construye a partir de esa voz--- tejiendo una tela inmortal.

¿Suena familiar? Por supuesto. Penélope y Circe son tejedoras. La simetría con que Homero nos presenta a estas dos formidables mujeres no puede estar más clara. Pero también hay una asimetría, como escribe Candelas Gala en su epílogo a \emph{Abandonando Ítaca}:

\begin{quotation}
Sí, porque Penélope y Circe pueden verse como dos caras de la misma moneda. Ambas son tejedoras, pero de un tipo muy diferente. [...] [Circe] era una gran conocedora de las hierbas y las pociones, una \emph{polifarmakos} según uno de sus epítetos homéricos. Su nombre evoca el círculo de Circe, que representa en su rueca cósmica. Homero la llamó Circe de las trenzas, sugiriendo en los nudos y trenzas de su cabello las fuerzas de creación y destrucción. Su tejido tenía el poder de gobernar las estrellas y los destinos de los hombres.

Penélope también era hilandera, enhebrando un sudario que tejía y destejía repetidamente. Mientras que el hilado de Circe era transformador y cambiante, el círculo vicioso de Penélope de repetición de lo mismo no podía competir con el tejido imprevisible y siempre renovable de Circe.
\end{quotation}

Pero no nos adelantemos a Homero. Uno de los aqueos, Polites, a quien Odiseo llama «el más atento y leal de los compañeros», propone que hagan ruido para manifestar su presencia. Circe los oye y los invita a pasar:

\begin{versocita}
Ella enseguida salió y, abriendo sus puertas brillantes,\\
los invitaba a pasar; iban todos tras ella, ignorantes…\\
Tan solo Euríloco fuera quedó, presintiendo un mal trance.\\
Ellos adentro, los hizo por tronos sentarse y sitiales,\\
y con cebada, con queso y miel verde aliñaba un brebaje en\\
vino de Pramnos, en el que ligó para el mismo potaje\\
drogas siniestras que hacían olvido del lar de los padres.\\
De que les dio la poción y bebieron, golpeó en ese instante\\
con su varita, y a las cochiqueras los lleva a guardarse.\\
Que de cochinos tenían la testa, la voz, la pelambre\\
y la figura, mas eran sus mientes las mismas de antes.
\end{versocita}

Fíjate en el detalle, biensabida lectora, siempreatento lector. Homero despacha la transformación en cerdos en unos pocos versos. Circe recurre a «drogas siniestras que hacían olvido del lar de los padres» y aquí encontramos otro de los interminables ecos que resuenan en la \emph{Odisea}\footnote{Como una moneda que va cayendo al fondo de un pozo, decía Pessoa.}. El filtro de Circe hace olvidar y evoca el de Helena, que impide sentir tristeza. Estamos ante poderosas hechiceras armadas con filtros ---y varitas mágicas---. ¿Hacía falta la elaborada escena de cirugía en la que la bruja transforma a los hombres de Ulises en gorrinos, como quien dice, «a mano»? ¿No podía el bienquerido director haber dedicado todos esos minutos a darle más fondo al personaje?

En la \emph{Odisea}, Euríloco no es transformado en cerdo, sino que escapa y le da noticias a su capitán.

\begin{versocita}
«Como ordenaste, cruzamos el bosque, preclaro Odiseo;\\
dimos por una honda hoz con casa hermosa, de lienzos\\
hechos de piedra pulida, y, en torno, un lugar bien abierto.\\
Alguien labraba al telar una tela cantando allí dentro,\\
una mujer o una diosa. A voces llamaron aquellos.\\
Ella enseguida salió, les abrió los portones espléndidos,\\
y los invita a pasar; todos iban detrás, bien ajenos…\\
Yo quedé fuera, que pálpito tuve de trance no bueno.\\
Todos, a un tiempo, invisibles se hicieron, y ni uno de ellos\\
apareció; que yo allí mucho rato me estuve en acecho».
\end{versocita}

Nada de chozas, nada de mujeres desconfiadas. Una casa hermosa de piedra pulida, una diosa aún más hermosa que canta. ¿Por qué insiste tanto Homero en la música de Circe? ¿Qué quiere decirnos? Ulises se dispone a explorar y le pide a su lugarteniente que le acompañe, pero este, muerto de miedo, le pide quedarse. Nuestro héroe, al que nadie puede negarle el arrojo, se dirige a casa de la bruja, cuando el dios Hermes se le cruza en el camino:

\begin{versocita}
Mas cuando, haciendo camino por la hoya encantada, me hallaba\\
ya de esa Circe milpocimera llegando a la casa,\\
Hermes varitadeoro, saliéndome al paso, me asalta\\
---la casa ya ante mis ojos--- igual que un muchacho en estampa\\
recién de barbiponiente, en la flor de la edad más graciada;\\
y me tomó de la mano, y así por mi nombre me hablaba:
\end{versocita}

Hermes no solo le pone sobre aviso, sino que le da los medios de enfrentarse a la hechicera:

\begin{versocita}
«¿Dónde, cuitado, otra vez marchas solo por estos rincones,\\
y sin saber lo que hay? Por la casa de Circe tus hombres\\
en cochiqueras cerradas guardados están, son cebones.
\stanzabreak
¿Vienes aquí a liberarlos? Te digo lo que aún no conoces:\\
téntalo y no volverás, quedarás donde están esos pobres.\\
Ah, pero vengo a salvarte y librarte de sus maldiciones.\\
Sí, vete a casa de Circe, mas este remedio antes tómate\\
benigno, que alejará de tu testa la hora sin nombre.
\stanzabreak
Voy a contarte el hacer pernicioso de Circe, y por su orden.\\
Preparará una mixtura y pondrá en la comida ese cóctel;\\
mas ni con eso hechizarte podrá, que a ese pisto se impone\\
la droga buena que ahora te doy; luego harás lo que oyes.
\stanzabreak
Cuando con larga varita intente ella darte un azote,\\
saca la espada aguzada que llevas al muslo y, de golpe,\\
a perseguirla, como deseando matarla, te pones.
\stanzabreak
Ella, temblando, a su cama invitarte querrá: nunca nones\\
debes a lecho de diosa decir si de cierto a tus hombres\\
los quieres ver liberados y que ella con mimo te aloje;\\
mas que te jure con gran juramento bendito de dioses\\
que ella no va a maquinar contra ti mal ninguno a mayores,\\
no se te lleve la hombría y valor cuando te desarropes».
\end{versocita}

Observa, detallista lectora, minucioso lector, que la aparición de Hermes no es casual. Odiseo es un simple mortal y es canónico que un mortal no puede enfrentarse solo a una diosa, lo que hace imprescindible que otro olímpico equilibre las tornas. El texto es delicioso. Frente al veneno, antídoto. Frente a la varita, espada. ¿Y qué hará la Circe cuando se vea acorralada? ¡Intentará encamar a Odiseo, por supuesto! Lo mejor de todo es que Hermes no le impone al héroe que se niegue, sino que le advierte de que antes de meterse en el tálamo con la divina, se asegure de que le restituya a sus hombres y le garantice hospitalidad. Nadie aquí parece preocupado por el hecho de que el valiente guerrero eche una cana al aire mientras Penélope espera, tejiendo y destejiendo su mortaja.

Odiseo llega a casa de Circe y empieza la acción:

\begin{versocita}
Ya en el portal de la trencigalana deidad, me detengo,\\
y allí parado le grito, y la diosa mi voz ya está oyendo.\\
Y ya sus puertas brillantes me abría saliendo en un vuelo\\
y me invitaba a pasar; la seguí con el ánimo inquieto.\\
Dentro, me lleva a sentar a un trono argenticlaveño\\
bello, de buena labor, con un lindo alzapié sobre el suelo.\\
Y me aliñó ella un brebaje en un cáliz de oro poniendo\\
en él la droga, tejiendo mi mal en su pecho malevo.\\
Y me lo dio, y lo bebí, pero no, no hubo allí encantamiento\\
por más que con su varita me diera, este ensalmo diciendo:
\stanzabreak
«Ve a la pocilga ya mismo, y acuéstate con tu pandilla».\\
Dijo, y, sacando yo entonces de al muslo mi espada buída,\\
me abalancé sobre Circe como por quitarle la vida.\\
Ella, gritando, se me agachó y se abrazó a mis rodillas\\
y entre sollozos palabras aladas así me decía:
\end{versocita}

\begin{versocita}
«¿Quién y de qué gentes tú? ¿Y tus padres? ¿En qué ciudad vives?\\
No entiendo cómo, aun bebida la pócima, embrujo no hice;\\
jamás un hombre ---jamás--- a esta droga logró resistirse\\
cuando, al beber, traspasó ella el redil de los dientes. Bebiste,\\
pero propósito hay dentro de ti que el hechizo no rinde.\\
Tú eres aquel milerrante Odiseo, el que siempre me dice\\
el Argifontes varitadeoro que habrá de venirme\\
un día aquí desde Troya en su léveda nave; y viniste.\\
Guarda esa espada en la vaina, y subamos tú y yo sin más quite\\
a nuestra cama, porque se nos venga que del confundirse en\\
lecho y amor tú confíes en mí y en ti yo confíe».
\end{versocita}

Tal y como había predicho Hermes. Ulises, que juega con ventaja, no se fía y le pide un juramento:

\begin{versocita}
Eso me dijo, y entonces hablé para así responderle:\\
«¿Cómo me puedes pedir que a ti, Circe, yo amable me acerque\\
si de mis hombres has hecho cochinos entre estas paredes\\
y a mí me tienes aquí y con artera cabeza me ofreces\\
entrar en tu dormitorio para que contigo me acueste y\\
puedas, cuando me desnude, la hombría quitarme y el temple?\\
No, no querría a tu cama subir, por más que te pese,\\
salvo que pases tú, diosa, por un juramento solemne\\
de que alumbrando no estás contra mí nuevo mal sobre este».
\end{versocita}

A la hechicera le falta tiempo para jurar, a Ulises para perdonarla y a Homero para meterlos en el catre:

\begin{versocita}
Dije, y juró ella enseguida, como se lo estaba pidiendo.\\
Y cuando ya hubo jurado y cerrado quedó el juramento,\\
con Circe entonces subí a su bienhermosísimo lecho.
\end{versocita}

¡Qué duros los trabajos del héroe! Mientras la comodona de su esposa se entretiene en la rueca, a él no le queda otra que darse unos revolcones divinos en el bienhermosísimo lecho de la hechicera.

Homero nos cuenta cómo las criadas de Circe lavan a Odiseo, lo perfuman, lo ungen con aceite y lo sientan a la mesa. Ulises, por una vez, parece desganado. La diosa le dice entonces:

\begin{versocita}
«¿A qué te sientas así, Odiseo, como un sin palabras\\
y pan y vino no tocas, mientras te recomes el alma?\\
¿Piensas que habrá un otro engaño? Pues de eso no debes en nada\\
desconfiar: me di a un juramento que es jura sagrada».
\end{versocita}

Lo cierto es que el héroe hace bien en no fiarse de los juramentos de los inmortales, tan aficionados a hacer promesas como a saltárselas. Además, Hermes le ha advertido que exija el regreso de sus hombres como prueba de buena voluntad.

\begin{versocita}
Dijo ella así, y respondiéndole yo le decía de vuelta:\\
«¿Ay, Circe, qué hombre iba a haber ---y digo hombre cabal--- que tuviera\\
de atiborrarse el valor, entregado a comienda y bebienda,\\
sin que libere a sus hombres primero y sus ojos los vean?\\
Que si me pides de corazón que yo coma y que beba,\\
suéltalos, y que a mis ojos mis buenos leales me vuelvan».
\end{versocita}

La hechicera no se hace de rogar, devuelve a sus compañeros la forma humana e invita a Ulises y a los suyos a arrastrar la nave hasta el arenal de la isla y a instalarse con ella. Ulises la obedece, regresa a la playa y le dice a su tripulación:

\begin{versocita}
«Al arenal arrastremos la nave primero y, ya en firme,\\
a los tesoros y jarcias daremos en cueva escondite;\\
luego aprestaos vosotros ---y no falte nadie--- a seguirme\\
porque a los vuestros veáis en la casa divina de Circe\\
bebiendo a bien y comiendo, que hay más de lo que un año rinde».
\end{versocita}

Euríloco, como era de esperar, no se fía y dice así:

\begin{versocita}
«Ay desgraciados, ¿ir dónde? ¿Por qué de más mal en amores?\\
¿Ir a la casa de Circe, que a todos nosotros a un golpe\\
en porcachones o en lobos nos convertirá o en leones,\\
para a la fuerza allí ser de esa su alta mansión guardadores?\\
Ved lo del Cíclope, cuando a su aprisco se entraron los hombres\\
---vuestros iguales--- con el temerario Odiseo a mayores:\\
la altanería de este de aquellos movió la hecatombe».
\end{versocita}

Tiene toda la razón, obviamente. ¿Y cuál es la reacción de Ulises?

\begin{versocita}
Dijo, y por mi pensamiento pasó, en esas dudando,\\
tirar del sable hojiluengo que va al recio muslo ajustado\\
y, tras cortar su cabeza, echarla a la tierra rodando,\\
y eso que me era pariente cercano; pero me frenaron\\
uno tras otro mis compañeros con estos halagos:
\stanzabreak
«Hijo del cielo, dejemos a este, si tú lo permites,\\
y aquí se quede junto a la nave y que nos la vigile;\\
y a los demás llévanos a esa santa morada de Circe».
\stanzabreak
Esto diciendo, dejaban el mar, el navío dejaban.\\
Pero ni Euríloco quiso quedarse en la nao concavada,\\
y nos seguía; temiendo a bien cierto mi atroz amenaza.
\end{versocita}

¡Un encanto este hombre! Euríloco es su lugarteniente, su amigo del alma y se ha librado por un pelo, nunca mejor dicho, de que le degollara por atreverse a llevarle la contraria. El caso es que vuelven a palacio y Circe les reitera a todos la invitación a pasarlo bien y descansar después de tantas cuitas. Los aqueos aceptan y se quedan... ¡por un año!
